\documentclass[../main.tex]{subfiles}

\begin{document}

\section{Introduction}\label{sec:intro}

This report addresses the design of an attitude control system for a launch
vehicle (LV) during the atmospheric phase of flight, at the instant of maximum
dynamic pressure ($t = 72$~s, max-$\bar q$). This is the hardest point to
control: the aerodynamic moment destabilises the airframe and the wind
disturbance is largest. I design the controller with classical
frequency-domain tools, using the Nichols chart as the main stability
indicator, following the short-period pitch-plane model of
Greensite~\cite{greensite1970}.

\subsection{Design methodology: frozen coefficients at max-$\bar q$}

Launch-vehicle attitude control in atmospheric flight is classically handled by
the \textbf{time-slice} (frozen-coefficient) method: the slowly varying mass,
inertia and aerodynamic coefficients are frozen at a representative instant, a
linear time-invariant short-period model is built and analysed at that operating
point, and the controller is then validated against time-varying
simulations~\cite[Sec.~2]{greensite_cr826}. The main sizing disturbance is
the wind profile, which drives both the trajectory dispersion and the structural
bending loads---hence the gust analysis of Tasks~1 and~2. I pick $t = 72$~s
because it is the load- and stability-critical slice: the aerodynamic
destabilising moment and the wind-induced angle of attack are both near their
maxima. The catch is that point-wise Nichols margins certify only the frozen
model and do not by themselves guarantee stability of the underlying
time-varying system; this is what motivates the parametric robustness study of
Task~3 and, eventually, gain scheduling along the trajectory.

\subsection{Plant model}

The pitch-plane short-period dynamics are described by the linear
time-invariant model (assignment Eq.~1), with state vector
$\mathbf{x} = [\,z,\ \dot z,\ \theta,\ \dot\theta,\ \eta,\ \dot\eta\,]^\top$
(lateral drift, pitch attitude and first bending mode), control input $\delta$
(TVC deflection) and disturbance $\alpha_w$ (wind angle of attack):
\begin{equation}
\dot{\mathbf{x}} =
\begin{bmatrix}
0 & 1 & 0 & 0 & 0 & 0\\
0 & a_1 & a_1 V + a_4 & 0 & 0 & 0\\
0 & 0 & 0 & 1 & 0 & 0\\
0 & A_6/V & A_6 & 0 & 0 & 0\\
0 & 0 & 0 & 0 & 0 & 1\\
0 & 0 & 0 & 0 & -\omega_1^2 & -2\zeta_1\omega_1
\end{bmatrix}\mathbf{x}
+
\begin{bmatrix} 0\\ a_3\\ 0\\ K_1\\ 0\\ -\phi_{1,TVC}T_c \end{bmatrix}\delta
+
\begin{bmatrix} 0\\ -a_1 V\\ 0\\ -A_6\\ 0\\ 0 \end{bmatrix}\alpha_w .
\label{eq:plant}
\end{equation}
The pitch channel is the core of the problem and follows from a moment balance
about the centre of mass. The aerodynamic moment $N_\alpha l_\alpha$ acts on the
local incidence---built from the pitch attitude $\theta$, the plunging term
$\dot z/V$ and the wind angle $\alpha_w$---while the gimballed thrust contributes
a control moment $T_c l_c\,\delta$; dividing the balance
$I_{yy}\ddot\theta = N_\alpha l_\alpha(\theta + \dot z/V - \alpha_w) + T_c l_c\,\delta$
by $I_{yy}$ reproduces row~4 of Eq.~\eqref{eq:plant},
\begin{equation}
\ddot\theta = A_6\,\theta + \frac{A_6}{V}\,\dot z + K_1\,\delta - A_6\,\alpha_w ,
\label{eq:pitch_channel}
\end{equation}
and identifies the two structural coefficients: the \textbf{aerodynamic moment}
$A_6 = N_\alpha l_\alpha / I_{yy} > 0$ (denoted $\mu_\alpha$) and the
\textbf{control effectiveness} $K_1 = T_c l_c / I_{yy}$ ($\mu_c$). With $A_6 > 0$
the airframe transfer function $\theta/\delta = K_1/(s^2 - A_6)$ carries an
open-loop pole at $s = +\sqrt{A_6} = +1.84$~rad/s in the right half-plane: the
\textbf{statically unstable} airframe of a vehicle whose centre of pressure lies
ahead of its centre of mass, for which feedback is
mandatory~\cite[Eqs.~(9.2.25)--(9.2.26)]{barrows2020}. The TVC coefficient $K_1$
sets the control authority available to counter it.

The sensed quantities are produced by the inertial unit (assignment Eq.~2),
which couples the bending coordinate $\eta$ into the gyro and accelerometer
channels through $\sigma_{1,ins}$ and $\phi_{1,ins}$:
\begin{equation}
\theta_m = \theta + \sigma_{1,ins}\,\eta,\quad
\dot\theta_m = \dot\theta + \sigma_{1,ins}\,\dot\eta,\quad
z_m = z - \phi_{1,ins}\,\eta,\quad
\dot z_m = \dot z - \phi_{1,ins}\,\dot\eta .
\label{eq:ins}
\end{equation}
This contamination is exactly what makes the lightly damped first bending
mode---written with $\omega_1 \equiv \omega_{BM} = 18.9$~rad/s and
$\zeta_1 \equiv \zeta_{BM} = 0.005$ in Eq.~\eqref{eq:plant}---visible to the
controller and drives the need for a dedicated filter (Task~2).

The actuator is the second-order TVC of Eq.~3,
\begin{equation}
W_{TVC}(s) = \frac{\omega_{TVC}^2}{s^2 + 2\zeta_{TVC}\omega_{TVC}s + \omega_{TVC}^2},
\qquad \omega_{TVC}=70~\text{rad/s},\ \zeta_{TVC}=0.7,
\label{eq:tvc}
\end{equation}
in series with a pure transport delay $\tau = 20$~ms. The delay is approximated
by a third-order Pad\'e rational function $e^{-\tau s} \approx P_3(s)$, the
diagonal $[3/3]$ member of the family obtained by matching the Taylor expansion
of $D(s)\,e^{-\tau s} = N(s)$ to the highest possible order; the illustrative
second-order form is
\begin{equation}
e^{-\tau s} \approx \frac{12 - 6\,\tau s + (\tau s)^2}{12 + 6\,\tau s + (\tau s)^2},
\label{eq:pade}
\end{equation}
an \textbf{all-pass} rational function ($|P(j\omega)| = 1$ for all $\omega$) that
reproduces the phase lag of the delay while leaving the loop gain unchanged.
I keep order three, matching the implementation, so that the approximation
stays accurate up to the bending frequency, where $\omega_{BM}\tau \approx 0.38$
is no longer negligible~\cite{franklin2015}. Parameters at $t=72$~s are read from the
reference data set \texttt{GreensiteLPV\_DATA.mat} (cross-checked against
Table~1 of the assignment); the salient values are reported in
Table~\ref{tab:params}.

\begin{table}[H]
\centering
\caption{Pitch-plane LV parameters at the max-$\bar q$ condition ($t=72$~s).}
\label{tab:params}
\begin{tabular}{lcc}
\toprule
Quantity & Symbol & Value \\
\midrule
Aerodynamic moment   & $A_6=\mu_\alpha$ & $3.382~\text{s}^{-2}$ \\
Control effectiveness & $K_1=\mu_c$     & $4.565~\text{s}^{-2}$ \\
Drift coefficients    & $a_1,a_3,a_4$   & $-0.0154,\ 20.61,\ -27.27~\text{s}^{-2}$ \\
Relative velocity      & $V$            & $937.7~\text{m/s}$ \\
Bending mode           & $\omega_{BM},\zeta_{BM}$ & $18.9~\text{rad/s},\ 0.005$ \\
INS coupling           & $\phi_{1,ins}$ [--], $\sigma_{1,ins}$ [rad/m] & $0.8,\ 0.178$ \\
TVC actuator           & $\omega_{TVC},\zeta_{TVC},\tau$ & $70~\text{rad/s},\ 0.7,\ 20~\text{ms}$ \\
\bottomrule
\end{tabular}
\end{table}

\subsection{Control architecture}

A proportional--derivative attitude law is adopted, augmented by a weak feedback
on the lateral-drift channel:
\begin{equation}
\delta_{cmd} =
K_{P,\theta}\,(\theta_{ref}-\theta_m) - K_{D,\theta}\,\dot\theta_m
- K_{P,z}\,z_m - K_{D,z}\,\dot z_m ,
\label{eq:control}
\end{equation}
with $K_{P,z},K_{D,z}$ small and negative, of order $10^{-3}$, as recommended in
the assignment guidelines. Closing the loop also on $(z_m,\dot z_m)$ trades
between two classical objectives in wind: containing the lateral
drift (drift minimum) and containing the angle of attack
$\alpha = \theta + \dot z/V + \alpha_w$, hence the structural load
$\bar q\,\alpha$ (load minimum)~\cite[Lecture~17]{dclvnotes} --- the binding concern at the
max-$\bar q$ condition studied here. The weak negative gains keep the drift
bounded without fighting the natural load-relieving weathercock of the
airframe; I therefore monitor both $z$ and $\bar q\,\alpha$ in the gust
simulations. The actual deflection follows from the actuator and
the bending filter $H_x(s)$ (Task~2): $\delta = W_{TVC}(s)\,e^{-\tau s}\,H_x(s)\,\delta_{cmd}$.

\subsection{A note on conditional stability}

Because the open-loop airframe carries a single right-half-plane pole
($P_R = 1$, the $+1.84$~rad/s aerodynamic pole), the loop is
\textbf{conditionally stable}. By the Nyquist criterion $N = P_R - Z_R$,
closed-loop stability ($Z_R = 0$) demands exactly one counter-clockwise
encirclement of the critical point, so stability is lost if the loop gain is
reduced below a threshold, not only if it is
increased~\cite[Eq.~(9.4.12), Figs.~9.17--9.18]{barrows2020}. An aerodynamically
unstable vehicle driven through a second-order (or higher) actuator therefore
has two gain margins---a low-frequency aerodynamic margin (reduce
gain $\to$ unstable) and a high-frequency rigid-body margin (increase
gain $\to$ unstable)---which on the Nichols chart show up as two $0$~dB crossings
of the $-180^\circ$ line. I therefore read the \texttt{margin} routine in
magnitude, and the assignment targets $|GM|\approx 6$~dB and
$|PM|\approx 30^\circ$ refer to this finite stable gain band. This is standard
practice for the attitude control of aerodynamically unstable
launchers~\cite{wie2008,franklin2015}.

\end{document}
